\documentclass[10pt,aspectratio=169,dvipsnames]{beamer}
\usetheme[]{Berlin}

\setbeamertemplate{footline}{
  \usebeamercolor[fg]{framesource}%
  \usebeamerfont{page number in head}%
  \hspace{0.2cm}
  \small \insertframenumber
  \vspace{0.2cm}
  \doclicenseIcon
  \hfill
  
\includegraphics[height=0.8cm]{logo-secondary-light.png}
  \hspace{0.2cm}
}

\setbeamercovered{invisible}

\setbeamertemplate{footline}[
 myframe number]

% PACKAGES
\usepackage[absolute,overlay]{textpos}
\usepackage[utf8]{inputenc}
\usepackage[official]{eurosym}
\usepackage{booktabs}
\usepackage{parskip}
\usepackage{bm}
\usepackage{tikz}
\usepackage{adjustbox}

\usepackage[
    type={CC},
    modifier={by},
    version={4.0},
]{doclicense}

% HYPERREFERENCES
\usepackage{hyperref}
\hypersetup{
	colorlinks=true,
	citecolor=tub-blue,
	linkcolor=tub-blue,
	urlcolor=tub-blue
}

% GRAPHICS
\graphicspath{
    {graphics/},
    {graphics/logos/},
  }
\DeclareGraphicsExtensions{.pdf,.jpeg,.png,.jpg}

% FORMATTING
\setlength{\parskip}{6pt}
\linespread{1.1}
\newcommand{\seprule}{\par\noindent\textcolor{black!25}{\rule{\textwidth}{0.4pt}}}

% SOURCES
\setbeamercolor{framesource}{fg=gray}
\setbeamerfont{framesource}{size=\tiny}
\newcommand{\source}[1]{\begin{textblock*}{5cm}(8.6cm,8.25cm)
    \begin{beamercolorbox}[ht=0.5cm,right]{framesource}
        \usebeamerfont{framesource}\usebeamercolor[fg]{framesource} Source: {#1}
    \end{beamercolorbox}
\end{textblock*}}

% REFERENCES
% \usepackage[style=nature, doi=true, maxcitenames=3]{biblatex}
% \bibliography{bibliography}
% \renewcommand*{\bibfont}{\footnotesize}

% COLOR TEXT
\newcommand{\bl}[1]{\textcolor{tub-blue}{#1}}
\newcommand{\gr}[1]{\textcolor{tub-green}{#1}}
\newcommand{\rd}[1]{\textcolor{tub-red}{#1}}
\newcommand{\yl}[1]{\textcolor{tub-yellow}{#1}}
\newcommand{\grey}[1]{\textcolor{gray}{#1}}
% \newcommand{\org}[1]{\textcolor{tub-orange}{#1}}
\newcommand{\cb}[1]{\colorbox{gray!20}{#1}}

\usepackage{array}% https://ctan.org/pkg/array
\makeatletter
\g@addto@macro{\endtabular}{\rowfont{}}% Clear row font
\makeatother
\newcommand{\rowfonttype}{}% Current row font
\newcommand{\rowfont}[1]{% Set current row font
   \gdef\rowfonttype{#1}#1%
}
\newcolumntype{L}{>{\rowfonttype}l}
\newcolumntype{R}{>{\rowfonttype}r}

\usepackage[normalem]{ulem}
\usepackage{multicol}


% FONT
\usepackage{roboto}

% REFERENCES
%\usepackage[backend=biber,style=authoryear-comp]{biblatex}
%\bibliography{references.bib}

% TITLE PAGE
\title{RESILIENT -- Resilient Energy System Infrastructure Layouts for Industry, E-Fuels and Network
Transitions}
\author{\url{https://resilient-project.github.io/}}

\institute[Technische Universität Berlin] % (optional, but mostly needed)
{
  \footnotesize
  Prof Dr. Tom Brown, \href{mailto:t.brown@tu-berlin.de}{t.brown@tu-berlin.de} \\
  Department of Digital Transformation in Energy Systems \\
  Technical University of Berlin, Germany \\~\\
  \alert{RESILIENT NRW Workshop}
}

\date{\today}

\subject{RESILIENT}

\titlegraphic{%
\vspace{-2.2cm}
\hspace{0.7cm}
\includegraphics[height=1.7cm,clip=true]{logo-secondary-light.png}
\hspace{0.7cm}
\includegraphics[trim=0 0cm 0 0cm,height=1.2cm,clip=true]{tub_logo.pdf}
\hspace{0.7cm}
\includegraphics[height=1.2cm,clip=true]{cetp_logo.png}
\hspace{0.7cm}\includegraphics[height=3cm,clip=true]{BMWE_Logo_mit_Förderzusatz_en.png}
}

\begin{document}

\addtocounter{framenumber}{-1}
{
	\setbeamertemplate{footline}{
		\usebeamercolor[fg]{framesource}%
		\usebeamerfont{page number in head}%
		\begin{center}
			\Large
			\doclicenseIcon
			\vspace{0.4cm}
		\end{center}
	}
	\maketitle
}

\begin{frame}{RESILIENT-Projektübersicht}
	\begin{itemize}
		\item \alert{Problemdiagnose}: Welt befindet sich in einer Phase der \alert{Unsicherheit} und des \alert{Wandels}
		      (Geopolitik, Energiewende, Klimawandel); Szenarien mit einzelnen deterministischen
		      Lösungen nicht mehr bedarfsgerecht
		\item \alert{Projektziel}: Entwicklung und Demonstration von Werkzeugen für eine resiliente Planung der Energiewende für Industrie, E-fuels und Netze
		\item Baut auf dem schon etablierten, quelloffenen und weltweit genutzten Multi-Vektor-Energiesystemmodell \alert{Python for Power System
			      Analysis (PyPSA)}
		\item \alert{Methodischer Schwerpunkt}: Optimierung unter Unsicherheit (stochastisch und robust),
		      Dekompositionsmethoden und Heuristiken für Resilienz
		\item \alert{Fallstudien}: \alert{Industriewende in NRW}, Netzplanung in BW, E-Fuels in
		      Skandinavien, Klimaextremereignisse in Frankreich
		\item Gefördert durch Clean Energy Transition (CET) Partnerschaft (EU-Kofinanzierung)
	\end{itemize}

\end{frame}

\begin{frame}{PyPSA-Eur: CO$_2$ and H$_2$ networks for self-sufficient net-zero Europe}

	\centering
	\includegraphics[height=\textheight]{carbon_networks_balance_map_hydrogen.png}

	\source{Hofmann, Tries, Neumann, Zeyen, Brown, 2024; \url{https://arxiv.org/abs/2402.19042}}

\end{frame}

\begin{frame}{Optimierung unter Unsicherheiten}

	\begin{columns}[T] % align columns
		\begin{column}{.5\textwidth}
			\begin{itemize}
				\item Unsicherheiten im Energiesystem erfordern Strategien, die in \alert{vielen Szenarien robust} sind
				\item Welche Investitionen in 2025-2030 erlauben die Erreichung von Klimazielen in allen Szenarien?
				\item Z.B. unter Verzögerungen beim Netzausbau, Knappheit von Wasserstoff, geopolitische Disruptionen
				\item Welche Maßnahmen \alert{schaffen Optionen} für künftiges Reagieren?
			\end{itemize}
		\end{column}
		\begin{column}{.5\textwidth}

			\includegraphics[width=\textwidth]{stochastic_optimisation.png}
		\end{column}
	\end{columns}
\end{frame}

\begin{frame}{Partners}
	\centering
	\includegraphics[width=0.9\textwidth]{resilient_partners.png}
\end{frame}

\begin{frame}{Work Plan}
	\centering
	\includegraphics[width=0.8\textwidth]{resilient_project_structure.png}
\end{frame}

\begin{frame}{Selection of Planned Model Developments}
	\begin{columns}[T] % align columns
		\begin{column}{.5\textwidth}
			\begin{minipage}[t][.45\textheight]{\linewidth}
				\begin{alertblock}{Computational Methods for Uncertainties}
					\begin{itemize}
						\item decomposition techniques
						\item large-scale stochastic optimisation
						\item test robustness of system
						\item using SMS++ framework
					\end{itemize}
				\end{alertblock}
			\end{minipage}

			\vspace{0.03\textheight} % separation between the boxes

			\begin{minipage}[t][.45\textheight]{\linewidth}
				\begin{alertblock}{Industry Transformation (FORECAST)}
					\begin{itemize}
						\item fuel and process switching
						\item industry relocation
						\item carbon sources and feedstocks
						\item data on stock \& investment cycles
						\item new technologies (oxyfuel cement, etc.)
					\end{itemize}
				\end{alertblock}
			\end{minipage}
		\end{column}

		\begin{column}{.5\textwidth}
			\begin{minipage}[t][.45\textheight]{\linewidth}
				\begin{alertblock}{Carbon Management and Biomass Usage}
					\begin{itemize}
						\item CO$_2$ network
						\item CO$_2$ sequestration potentials
						\item circular carbon economy and recycling
						\item biomass usage options
					\end{itemize}
				\end{alertblock}
			\end{minipage}

			\vspace{0.03\textheight} % separation between the boxes

			\begin{minipage}[t][.45\textheight]{\linewidth}
				\begin{alertblock}{Global Green Fuel and Material Markets}
					\begin{itemize}
						\item imports of green energy and materials
						\item effects on European infrastructure
						\item restructuring of value chains
						\item risks (geopolitical, technological, etc.)
					\end{itemize}
				\end{alertblock}
			\end{minipage}
		\end{column}
	\end{columns}
\end{frame}

\begin{frame}{Fallstudie Industrietransformation in Nordrhein-Westfalen}

	Warum NRW? Viel Industrie, Einbettung Europa, Strategien für CO$_2$ \& H$_2$(-Importe)

	\centering
	\includegraphics[height=0.85\textheight]{sequestration_map.png}

	\source{\href{https://www.energy4climate.nrw/industrie-produktion/co2-management/umgang-mit-co2}{NRW.Energy4Climate}}
\end{frame}

\begin{frame}{Fallstudie Industrietransformation in Nordrhein-Westfalen}
	\begin{itemize}
		\item \alert{Problemdiagnose}: Transformation der energieintensiven Industrien erfordert Umstellung von sowohl Brennstoffen als auch Prozessen, Investitionszyklen sind lang, abhängig von verfügbaren Infrastrukturen, viele Unsicherheiten
		\item \alert{RESILIENT-Ansatz}: Herausarbeitung und Konkretisierung der Probleme mit Partnern aus der Industrie $\Rightarrow$ \alert{heutiger Workshop}; Entwicklung von Strategien, die über mehrere Szenarien robust sind
		\item Beispiele für \alert{Unsicherheiten}: Preise und Verfügbarkeit von grünem Wasserstoff; Verfügbarkeit CO$_2$-Pipelines und Biomasse; politische Unterstützung und Marktdynamiken (KSV, grüne Leitmärkte); Akzeptanz für Infrastruktur, CCS, usw.
		\item Beispiele für \alert{Strategien}: Hybridansätze, z.B. bivalente Kessel, Nutzung von (Bio-)Methan für DRI bis H$_2$ da ist; Elektrolyseure vor Ort, CO$_2$-Transport per Schiff/Zug/Lkw, Speicherung von Strom/H$_2$/Derivaten, Flexibilisierung, blauer H$_2$, Import von Zwischenprodukten (HBI, Ammoniak, Methanol)
	\end{itemize}

\end{frame}

\begin{frame}{Ziel des Workshops heute}
	\begin{itemize}
		\item Verständnis für die Unsicherheiten und Hedging-Strategien schaffen
		\item Wo sind die \alert{entscheidenden Herausforderungen}?
		\item Welche Technologien können Robustheit gewährleisten?
		\item Wie entwickeln wir unsere Fallstudie, sodass sie für die Praxis relevant bleibt?
	\end{itemize}

\end{frame}

\end{document}
